\documentclass[11pt,a4paper]{article}
\usepackage[utf8]{inputenc}
\usepackage[T1]{fontenc}
\usepackage[french]{babel}
\usepackage{amsmath,amssymb,booktabs,hyperref}
\title{Fiche de lecture — Projet 3}
\author{Groupe 3 — ENSA Safi}
\date{Semaine 1 — Phase 0}
\begin{document}
\maketitle

\section{Objet de la lecture}
Le projet étudie l'information imparfaite, les équilibres bayésiens parfaits
(EBP/PBE) et les mécanismes d'enchères. Cette fiche fixe les notions et les
formules à valider avant l'implémentation Python.

\section{Équilibre bayésien parfait}
Un EBP est un profil de stratégies $\sigma$ et de croyances $\mu$ satisfaisant
la rationalité séquentielle et la cohérence des croyances. Pour un message $m$
atteint avec une probabilité positive :
\[
\mu(t\mid m)=\frac{\pi(t)\Pr(m\mid t)}
{\sum_{t'}\pi(t')\Pr(m\mid t')}.
\]
La stratégie du récepteur est une meilleure réponse :
\[
a^*(m)\in\arg\max_{a\in A}\sum_t\mu(t\mid m)u_R(t,m,a).
\]
Chaque type de l'émetteur ne doit pas gagner à dévier :
\[
\sigma(t)\in\arg\max_{m\in M}u_E(t,m,a^*(m)).
\]
Pour un message hors chemin, la croyance doit être une distribution
normalisée et la rationalité séquentielle reste exigée.

\section{Enchères}
Pour des valeurs privées indépendantes $v_i\sim U[0,1]$, l'enchère symétrique
de premier prix sert de repère :
\[
b(v)=\frac{n-1}{n}v.
\]
Dans une enchère au second prix, la stratégie dominante est l'offre sincère
$b(v)=v$. Le revenu et les gains seront estimés par simulation Monte-Carlo :
\[
\widehat{\mathbb E}[X]=\frac{1}{N}\sum_{k=1}^{N}X_k,\qquad
IC_{95\%}\simeq \widehat{\mathbb E}[X]\pm1.96\frac{s_X}{\sqrt N}.
\]

\section{Approche de calcul prévue}
\begin{enumerate}
  \item Formaliser les types, croyances, messages, actions et utilités.
  \item Énumérer les profils de signalisation candidats (séparants et
        regroupants).
  \item Calculer les croyances par Bayes sur le chemin et tester les croyances
        hors chemin.
  \item Vérifier la meilleure réponse du récepteur puis les déviations de
        chaque type.
  \item Simuler les enchères avec une graine reproductible et comparer aux
        repères théoriques.
\end{enumerate}

\section{Ressources du cours}
\begin{itemize}
  \item \href{https://www.gametheory.net/}{Game Theory — ressources de référence}.
  \item \href{https://www.cs.cmu.edu/~epxing/Class/10708-19/notes/lecture-07.pdf}
        {Notes sur l'information incomplète et les croyances}.
  \item \href{https://ollama.com/library/llama3.2}{Ollama — llama3.2}.
\end{itemize}
\textit{Les liens et conventions seront complétés après validation des supports
officiels du module.}
\end{document}
